\section{The common orientation source}
\label{sec:orientation-source}

This section keeps track of the sign lost by projectivizing the Clebsch
chart.  It first normalizes the pulled-back incidence cover, then follows
deck exchange through the golden six-axis data and the Petersen coefficient
module.

\subsection{Normalization over the Clebsch chart}

Let \(\mathcal N\to\mathbf P(H)\) be the incidence normalization from
Section~\ref{sec:cover}, and put \(B=\mathbf P(V_E)\).  On the chart,
Theorem~\ref{thm:arithmetic-main} and Hitchin's restriction identity give
\[
 z^2=5\iota_t^*J_0=80\sigma_3^2
   =(z-4\sqrt5\,\sigma_3)(z+4\sqrt5\,\sigma_3).
\]
Each factor has function field \(E(B)\), and \(B\) is normal.  The
normalization of \(B\times_{\mathbf P(H)}\mathcal N_E\) is therefore two
copies \(B_+\amalg B_-\).  The constant plane \(U_t\) supplies one
component; deck exchange supplies the other.  On \(D(\sigma_3)\), the two
linear factors differ by a unit, so the original pullback is already their
product and is finite etale.  Hitchin's chart classification identifies
this as exactly the locus of two distinct regular configurations.  Along
\(\sigma_3=0\), the factors meet before normalization, but their normal
components remain separate.

This factorization is an equality of sections after the same local
trivialization of \(\mathcal O(3)\) used in the square-class proof.  It does
not strengthen the integral statement: comparison with the geometric
incidence model still holds only over an unspecified \(\mathbf Z[1/N]\).

\subsection{The golden involution}

For the ordered representatives \(a_0,\ldots,a_5\) of \(I_t\), their Gram
matrix has the form
\[
 G=(t+2)I+tC,
\qquad
C=\begin{pmatrix}
0&1&1&1&-1&-1\\
1&0&-1&-1&-1&-1\\
1&-1&0&1&1&-1\\
1&-1&1&0&-1&1\\
-1&-1&1&-1&0&-1\\
-1&-1&-1&1&-1&0
\end{pmatrix}.
\]
Direct summation of the six rank-one matrices gives
\(\sum_i a_ia_i^{\mathsf T}=2(t+2)I_3\), hence
\(G^2=2(t+2)G\).  Substituting \(G=(t+2)I+tC\) and using
\((t+2)^2=5t^2\) gives \(C^2=5I\).  Changing signs of axis
representatives switches \(C\) to \(DCD\), but every triangle product and
hence \(Z_C\) is unchanged.  The off-diagonal entries of \(C^2\) also give
\[
 \sum_{k\ne i,j}C_{ij}C_{jk}C_{ki}=0.
\]
Thus the directional derivative of \(Z_C\) along \(\mathbf1\) vanishes, so
\(Z_C(x+u\mathbf1)=Z_C(x)\) and the cubic descends to
\(\Q^6/\Q\mathbf1\).

Let \(a_i'\) denote the conjugate axes, obtained by \(t\mapsto1-t\), and
let \(R\) be the exchanger of Section~\ref{sec:arithmetic}.  Direct
substitution gives
\[
 Ra_i=t d_i a'_{p(i)},\qquad
 p=(0,1,5,4,2,3),\qquad
 (d_i)=(1,-1,1,1,1,1).
\]
Using \(t^2(1-t)=-t\), one obtains
\[
 d_id_jC_{p(i)p(j)}=-C_{ij}.
\]
Thus, at the golden fibre, conjugation transported by \(R\) changes \(C\)
to \(-C\) and \(Z_C\) to \(-Z_C\).  The lone negative \(d_i\) is only a change of
axis representatives and does not cause the reversal; switching alone
leaves the triangle cubic fixed.

No global marking of the varying configurations on \(B_-\) is needed or
asserted.  The fixed section \(U_t\) marks \(B_+\) by the constant source
\([C,Z_C]_{\rm or}\).  We mark \(B_-\) by its deck-opposite, meaning that
the normalized odd generator changes sign.  The calculation at \([xyz]\)
proves that this algebraic convention agrees with the actual conjugate
golden configuration there; it is not used to infer a pointwise conference
labeling elsewhere.  Since the normalized components remain disjoint above
\(\sigma_3=0\), the two source labels extend there even though the literal
two-configuration interpretation fails.

\subsection{Transport to degree six}

Write Hitchin's chart as
\(\iota_t(y)=\sum_i y_iq_i^{(t)}\), with \(q_1=xyz\).  Since
\(R(xyz)=-xyz\), equivariance and irreducibility of the four-module give
\[
 Rq_i^{(t)}=-q_i^{(1-t)}
\]
after transporting the five plane-triple labels.  Normalize the lift on
\(B_+\) by \(q_1=xyz\); the deck-opposite source carries the lift
\(-\iota_t\).  The displayed exchanger calculation shows that this
normalization agrees with golden conjugation at \([xyz]\).  Thus the lift
is part of the oriented source datum, rather than a purported global marking
of the second family.

The primitive integral map
\[
 \beta(y)_{ij}=y_i+y_j
\]
lands in the Petersen \((-2)\)-eigenspace and satisfies
\[
 \sum_{i<j}(y_i+y_j)^2=3\sum_i y_i^2,
 \qquad
 y_i=\frac13\sum_{j\ne i}\beta(y)_{ij}.
\]
Multiplicity one makes this the canonical comparison after the five labels
are fixed.  Reversing the chart lift reverses \(y\), the zonal harmonic
\(F_y\), and its cubic integral.  The two signs in the normalized incidence
cover therefore select the two signs of the exact Gaunt cubic.  This proves
Theorem~\ref{thm:orientation-source}.

There is no ambient linear comparison between harmonic degrees three and
six: their irreducible rotation modules have dimensions seven and thirteen.
The comparison exists only on the chosen Clebsch four-space, and it stops at
the oriented source class \([C,Z_C]_{\rm or}\).
